% =====================================================================
%  CAPÍTULO 6 -- RESULTADOS E CONSIDERAÇÕES FINAIS
%  (Regulamento, Art. 9º, item vi; Manual, seção 2.2.3)
%  O Manual recomenda que a parte conclusiva NÃO contenha citações.
% =====================================================================

\chapter{Resultados e considerações finais}
\label{cap:resultados}

Este projeto conduziu a concepção do RoundUp como um processo de
ideação, no qual as ideias de funcionalidades foram convertidas em
requisitos e esses requisitos foram usados, ao mesmo tempo, para
construir o aplicativo e para avaliar as ideias com usuários reais. Essa estratégia permitiu rastrear cada
funcionalidade desde a ideia até a decisão de mantê-la, ajustá-la ou
revisá-la, e daí até o código e os testes.

\section{Resultados alcançados}

Os objetivos específicos apresentados no \autoref{cap:introducao} foram
atingidos da seguinte forma:

\begin{enumerate}
  \item \textbf{Fatores de motivação e evasão.} Atingido de forma
        exploratória e limitada. Os praticantes reais indicaram a saúde
        mental, a saúde física e o convívio com a turma como motivos para
        continuar treinando, e os dois ex-praticantes pararam por falta de
        tempo ou conflito com a rotina. Ambos afirmaram que um aplicativo
        não teria mudado sua decisão, o que sugere que parte da evasão tem
        causa externa ao treino. A visão dos professores sobre a evasão não
        pôde ser levantada.
  \item \textbf{Ideias convertidas em requisitos rastreáveis.} Atingido.
        Treze ideias originaram, junto com os requisitos de base, 33
        requisitos funcionais, 9 regras de negócio e 11 requisitos não
        funcionais, todos registrados com origem, ciclo, prioridade e
        situação.
  \item \textbf{Avaliação com usuários reais.} Atingido, com amostra
        pequena. Seis respondentes avaliaram as ideias; os recursos de
        progresso visível tiveram as melhores notas, e lembretes e
        mensagens automáticas, as piores. Os comentários reforçaram
        requisitos de baixo esforço de registro, de controle de
        notificações e de privacidade.
  \item \textbf{Personas e teste do instrumento.} Atingido. Doze personas
        sintéticas orientaram a geração de 65 respostas de teste, enviadas
        a uma cópia do formulário. O teste confirmou o roteamento das três
        seções, inclusive a do professor, e o funcionamento do procedimento
        de análise, que fica pronto para uma nova coleta real. Como insumo
        complementar da ideação, as personas levaram o autor a revisitar
        os requisitos, principalmente os voltados a professores e
        academias; as respostas sintéticas foram mantidas separadas dos
        dados reais e não entraram nas notas que definiram a situação de
        cada requisito.
  \item \textbf{Modelagem e implementação do protótipo.} Atingido. O
        protótipo é um aplicativo funcional com três perfis de usuário. Os 33
        requisitos funcionais foram implementados, incluindo a data de
        conquista das medalhas, o público-alvo e a vigência das missões, a
        personalização das faixas e o envio de notificações, verificados
        em um iPhone. As regras de negócio que protegem a gamificação,
        como a validação humana antes do crédito de pontos, foram
        implementadas nas regras de segurança do banco de dados e
        cobertas por 288 testes automatizados.
  \item \textbf{Conformidade com a LGPD.} Atingido. O consentimento no
        cadastro, o registro do responsável legal de alunos menores de
        idade, a guarda dos dados pessoais em documento visível só ao
        próprio usuário, o isolamento dos dados entre academias e professores e a
        exclusão de conta com remoção dos dados em cascata foram
        implementados. Fica em aberto se o registro do consentimento deve
        ser preservado após a exclusão, como prova, o que depende de
        orientação jurídica.
\end{enumerate}

\section{Considerações sobre o uso de respostas sintéticas}

As personas e as respostas sintéticas foram úteis para dois fins: como
insumo complementar da ideação, ao dar voz a perfis que a coleta real não
alcançou e levar o autor a repensar os requisitos, e como dados de teste
do questionário e da análise, principalmente na seção do professor. Elas não substituem a coleta com
usuários reais, porque reproduzem padrões supostos por quem as gerou;
quando ficam próximas dos dados reais, isso mostra apenas que o gerador
copiou bem o padrão que lhe deu origem. A lição prática é manter os dois
conjuntos sempre separados, identificar cada resposta sintética e
registrar o método de geração, para que nenhuma conclusão sobre os
usuários dependa delas.

\section{Limitações}

\begin{itemize}
  \item A amostra real é pequena (seis respondentes), pouco diversa
        (todos de Jiu-Jitsu) e não incluiu professores; seus resultados
        servem para orientar decisões, não para generalizar.
  \item Nenhum professor real respondeu, de modo que os requisitos de
        gestão não foram avaliados; as respostas sintéticas de professor
        servem apenas como dados de teste.
  \item Os valores-padrão de gamificação (pontos por presença, níveis e
        marcos de medalha) não foram calibrados com uso real, e o lembrete
        diário agendado das notificações foi testado apenas em emulador.
  \item O efeito do aplicativo sobre a evasão não pôde ser medido no prazo
        do TCC, e parte da evasão tem causas externas, como falta de tempo
        e lesões, que estão fora do alcance de um aplicativo.
\end{itemize}

\section{Trabalhos futuros}

\begin{itemize}
      \item Validar em aparelho o lembrete diário agendado e fazer o toque
        na notificação levar à tela correspondente (RF31).
  \item Permitir o uso sem conexão (RNF06) nas ações que não dependem de
        validação imediata do servidor, mantendo o \textit{check-in} e a
        validação de presença dependentes de conexão.
  \item Reaplicar o questionário com uma amostra maior, incluindo
        outras modalidades, concentrando a investigação nas ideias marcadas
        como ``a revisar''.
  \item Aplicar o questionário a professores e donos de academia reais,
        para avaliar os requisitos de gestão hoje marcados como ``sem
        dados''.
  \item Realizar um estudo de acompanhamento, medindo a frequência dos
        alunos antes e depois do uso do aplicativo.
\end{itemize}
